\section*{Chapter \ref{ch:b2:chromatography} --- Chromatography}

\begin{solution}{exo:b2:chromatography:1}
$k = (5.0 - 1.0)/1.0 = 4.0$; fraction of time in the \omterm{def:b2:chromatography:principle}{mobile phase} $1/(1 + k) = 0.20$.
\end{solution}

\begin{solution}{exo:b2:chromatography:2}
$N = 16(6.0/0.30)^2 = 16 \times 400 = 6400$.
\end{solution}

\begin{solution}{exo:b2:chromatography:3}
$H = \qty{250}{mm}/10\,000 = \qty{0.025}{mm} = \qty{25}{\micro m}$.
\end{solution}

\begin{solution}{exo:b2:chromatography:4}
Ethanol in blood: GC (volatile). Protein: HPLC. Caffeine in a drink: HPLC (in water, not volatile
enough without preparation). Benzene in petrol: GC. Sugars: HPLC (they decompose before they boil).
\end{solution}

\begin{solution}{exo:b2:chromatography:5}
$R_s = 2(4.6 - 4.0)/(0.40 + 0.50) = 1.2/0.90 \approx 1.3$: not quite baseline separated (below 1.5).
\end{solution}

\begin{solution}{exo:b2:chromatography:6}
$\sqrt N = 4R_s\,\dfrac{\alpha}{\alpha - 1}\,\dfrac{1 + k_2}{k_2} = 4 \times 1.5 \times 11 \times 1.25 =
82.5$, so $N \approx 6.8 \times 10^3$.
\end{solution}

\begin{solution}{exo:b2:chromatography:7}
$u_{\mathrm{opt}} = \sqrt{8/2} = \qty{2}{mm/s}$; $H_{\min} = 5 + 2\sqrt{16} = \qty{13}{\micro m}$.
\end{solution}

\begin{solution}{exo:b2:chromatography:8}
Phenol (polar \ce{OH}, hydrogen bonds with water) first, then benzene, then toluene (one more
\ce{CH3}, more nonpolar). More methanol makes the \omterm{def:b2:chromatography:principle}{mobile phase} less polar: all three elute earlier, in
the same order.
\end{solution}

\begin{solution}{exo:b2:chromatography:9}
$F = (860/400)/(20.0/10.0) = 2.15/2.00 = 1.075$. Sample: $c_X = (645/410) \times 10.0/1.075 \approx
\qty{14.6}{mg/L}$.
\end{solution}

\begin{solution}{exo:b2:chromatography:10}
$R_s \propto \sqrt N \propto \sqrt L$: $L$ must be multiplied by $(1.5/1.06)^2 \approx 2.0$, which doubles
the analysis time and the pressure. Other levers: the selectivity (\omterm{def:b2:chromatography:principle}{mobile phase}, \omterm{def:b2:chromatography:principle}{stationary phase},
temperature), smaller particles, the optimum flow rate.
\end{solution}

\begin{solution}{exo:b2:chromatography:11}
The peaks are $4\sigma R_s$ apart, so the midpoint is $2\sigma R_s$ from each. $R_s = 1$: each peak
contributes $\mathrm e^{-(2)^2/2} = \mathrm e^{-2} = 0.135$; the valley is at $0.27$ of the peak height.
$R_s = 1.5$: $2\mathrm e^{-(3)^2/2} = 2\mathrm e^{-4.5} \approx 0.022$, about 2~\%: back to the baseline
for practical purposes.
\end{solution}

\begin{solution}{exo:b2:chromatography:12}
$k_2/(1+k_2)$: 0.67, 0.83, 0.91; time in units of $t_M$: 3, 6, 11. Going from 2 to 5 gains 25~\% in
\omterm{def:b2:chromatography:separation}{resolution} for twice the time; from 5 to 10, 9~\% for nearly twice again. A $k$ from about 2 to 5 is
the usual compromise.
\end{solution}

\begin{solution}{pb:b2:chromatography:1}
\textbf{1.} Stationary: silica grains carrying long alkyl chains, nonpolar. Mobile: water with
methanol, polar.
\textbf{2.} They are very polar and stay in the \omterm{def:b2:chromatography:principle}{mobile phase}: $k \approx 0$.
\textbf{3.} Caffeine has one more methyl group than theophylline, and theophylline an \ce{N-H} that
hydrogen-bonds with water: theophylline is more polar and elutes first.
\textbf{4.} It is close in structure (similar response and behaviour), absent from the drink, and elutes
near caffeine while being resolved from it.
\textbf{5.} Both would decrease: a less polar \omterm{def:b2:chromatography:principle}{mobile phase} competes better with the \omterm{def:b2:chromatography:principle}{stationary phase}.
\textbf{6.} The time the \omterm{def:b2:chromatography:principle}{mobile phase} takes to cross the column, the time a compound spends moving.
\textbf{7.} $k_{\mathrm{theo}} = (3.0 - 1.0)/1.0 = 2.0$; $k_{\mathrm{caf}} = 2.4$.
\textbf{8.} $\alpha = 2.4/2.0 = 1.2$.
\textbf{9.} Caffeine $N = 16(3.4/0.20)^2 = 4624 \approx 4.6 \times 10^3$; theophylline $16(3.0/0.18)^2
\approx 4.4 \times 10^3$.
\textbf{10.} $H = \qty{150}{mm}/4624 \approx \qty{0.032}{mm} = \qty{32}{\micro m}$.
\textbf{11.} $R_s = 2(3.4 - 3.0)/(0.18 + 0.20) = 0.80/0.38 \approx 2.1$.
\textbf{12.} $\frac{\sqrt{4624}}{4} \times \frac{0.2}{1.2} \times \frac{2.4}{3.4} = 17 \times 0.167 \times
0.706 \approx 2.0$; the small difference comes from the equal-width assumption of the equation.
\textbf{13.} Yes: $R_s > 1.5$.
\textbf{14.} $N$ is halved: $R_s \approx 2.1/\sqrt2 \approx 1.5$, just at baseline separation.
\textbf{15.} Halved: caffeine at \qty{1.7}{min}.
\textbf{16.} $R_s \approx 2.1\sqrt2 \approx 3.0$, for twice the time and twice the pressure: a waste here.
\textbf{17.} The \omterm{def:b2:chromatography:principle}{mobile phase} composition (methanol fraction, pH, another organic solvent) or the
\omterm{def:b2:chromatography:principle}{stationary phase}.
\textbf{18.} Little: $k_2/(1 + k_2) = 0.71$ already; at $k_2 = 5$ it would be 0.83 (+18~\%) for an
analysis time multiplied by $6/3.4 \approx 1.8$.
\textbf{19.} $F = (1500/1000)/(50.0/40.0) = 1.50/1.25 = 1.20$.
\textbf{20.} Both peaks come from the same injection: the volume cancels in the ratio of areas.
\textbf{21.} $c = (970/1010) \times 40.0/1.20 \approx \qty{32.0}{mg/L}$.
\textbf{22.} Tenfold dilution: \qty{320}{mg/L}.
\textbf{23.} $A_{IS}$ would be too large, and the caffeine result too low.
\textbf{24.} A series of caffeine standards injected under exactly the same conditions as the sample, a
calibration line of area against concentration, and reproducible injection volumes.
\textbf{25.} $\qty{320}{mg/L} \times \qty{0.250}{L}$: $\boldsymbol{\approx \qty{80}{mg}}$ of caffeine per can (exercise
data).
\end{solution}
